\documentclass[11pt,a4paper]{article}
\PassOptionsToPackage{obeyspaces,spaces}{url}

\usepackage[margin=2.2cm]{geometry}
\usepackage[T1]{fontenc}
\usepackage{lmodern}
\usepackage{microtype}
\usepackage{xcolor}
\usepackage{booktabs}
\usepackage{tabularx}
\usepackage{enumitem}
\usepackage{listings}
\usepackage{amssymb}
\usepackage[most]{tcolorbox}
\usepackage{hyperref}
\usepackage{xurl}

\definecolor{accent}{HTML}{1F4E79}
\definecolor{codebg}{HTML}{F4F6F8}
\definecolor{good}{HTML}{2E7D32}
\definecolor{warn}{HTML}{B26A00}
\definecolor{bad}{HTML}{C62828}

\hypersetup{colorlinks=true, linkcolor=accent, urlcolor=accent}
\setlist{itemsep=2pt, topsep=4pt}
\setlength{\parindent}{0pt}
\setlength{\parskip}{5pt}

\lstset{
  language=bash,
  basicstyle=\ttfamily\small,
  backgroundcolor=\color{codebg},
  frame=single, rulecolor=\color{codebg},
  breaklines=true, columns=fullflexible, keepspaces=true,
  xleftmargin=4pt, xrightmargin=4pt,
  showstringspaces=false,
  commentstyle=\color{gray}\itshape,
  keywordstyle={}, morecomment=[l]{\#}
}

\newcommand{\p}[1]{\nolinkurl{#1}}  % paths / commands inline, with line breaks allowed

\newtcolorbox{step}[1]{colback=white, colframe=accent, fonttitle=\bfseries, title=#1,
  arc=2pt, boxrule=0.8pt, breakable, before skip=10pt, after skip=10pt}
\newtcolorbox{expect}{colback=good!6, colframe=good, fonttitle=\bfseries, title=What you should see,
  arc=2pt, boxrule=0.5pt, breakable}
\newtcolorbox{careful}[1][Careful]{colback=warn!6, colframe=warn, fonttitle=\bfseries, title=#1,
  arc=2pt, boxrule=0.5pt, breakable}

\title{\textbf{How to Run the Pacman Contest}\\[4pt]\large An exact, step-by-step guide for this MacBook}
\author{Team \texttt{pacman-agent}}
\date{October 2026}

\begin{document}
\maketitle

\begin{tcolorbox}[colback=accent!5, colframe=accent, title=\textbf{The 30-second version}, arc=2pt]
\begin{lstlisting}
cd /Users/bortekstor/Desktop/Eutopia_Pacman_Contest/Pack_It_Up/pacman-agent
source venv/bin/activate
cd pacman-contest/src/contest
python capture.py -r baseline_team -b ../../../my_team.py
\end{lstlisting}
This is the command from the organisers' \texttt{README.md}. Everything below explains each step exactly.
\end{tcolorbox}

\tableofcontents

%=====================================================================
\section{How the files are organised}
%=====================================================================

\begin{center}
\begin{tabularx}{\linewidth}{@{}lX@{}}
\toprule
\textbf{Path (inside \texttt{pacman-agent/})} & \textbf{What it is} \\
\midrule
\p{my_team.py} & \textbf{Your team. The only file you edit and submit.} It stays here, at the root. \\
\p{TEAM.md} & Your team details (fill in before submitting). \\
\p{venv/} & The Python environment (Homebrew Python 3.14 with a working Tk). \\
\p{pacman-contest/} & The game engine from the organisers. Do not edit. \\
\p{pacman-contest/src/contest/} & Where \texttt{capture.py}, \texttt{baseline\_team.py} and \texttt{layouts/} live. You run games from here. \\
\p{pacman-contest/agents/} & Old example teams from the engine repository. \textbf{Not} where your file goes. \\
\bottomrule
\end{tabularx}
\end{center}

\begin{careful}[Why \texttt{my\_team.py} is not in \texttt{agents/}]
\begin{itemize}
  \item The organisers' \texttt{README.md} says: \emph{``In the root folder \ldots\ create in \texttt{my\_team.py} a class \ldots''}
        and runs it as \p{../../../my_team.py}, which is the root of \texttt{pacman-agent}.
  \item \p{pacman-contest/agents/team_template/my_team.py} is an outdated example. It uses \texttt{import util}
        instead of \texttt{import contest.util}, so it does not load with the current engine.
  \item In the tournament, the organisers copy your \texttt{my\_team.py} into their own \p{teams/<your_team>/} folder.
        You don't create that folder yourself.
\end{itemize}
\end{careful}

%=====================================================================
\section{Step-by-step: your first game}
%=====================================================================

\begin{step}{Step 1 --- Open Terminal and go to the project folder}
Open \textbf{Terminal} (Cmd+Space, type \texttt{Terminal}, Enter). Then type:
\begin{lstlisting}
cd /Users/bortekstor/Desktop/Eutopia_Pacman_Contest/Pack_It_Up/pacman-agent
\end{lstlisting}
\end{step}

\begin{step}{Step 2 --- Activate the virtual environment}
\begin{lstlisting}
source venv/bin/activate
\end{lstlisting}
\begin{expect}
Your prompt now starts with \texttt{(venv)}. Check that the right Python is used:
\begin{lstlisting}
python -c "import tkinter, contest; print(tkinter.TkVersion)"
\end{lstlisting}
It must print \texttt{9.1}. If it prints \texttt{8.5} or an error, see Section~\ref{sec:trouble}.
\end{expect}
Repeat Step~2 in \textbf{every new Terminal window}. To leave the environment later, type \texttt{deactivate}.
\end{step}

\begin{step}{Step 3 --- Go to the engine folder}
\begin{lstlisting}
cd pacman-contest/src/contest
\end{lstlisting}
All game commands below are run from \textbf{this} folder. Here the short names \texttt{baseline\_team} and
\texttt{officeCapture} work, and your team is \p{../../../my_team.py}.
\end{step}

\begin{step}{Step 4 --- Watch a game}
\begin{lstlisting}
python capture.py -r baseline_team -b ../../../my_team.py
\end{lstlisting}
\begin{expect}
A game window opens. \textbf{Red (left)} is \texttt{baseline\_team}, \textbf{blue (right)} is your \texttt{my\_team.py}.
Each agent gets 300 moves. At the end Terminal prints the result, e.g.\ \texttt{Tie game!} or
\texttt{The Blue team wins by 5 points.}, and the window closes.
\end{expect}
To stop early: click on Terminal and press \textbf{Ctrl+C}.
\end{step}

\begin{careful}[Your \texttt{my\_team.py} is still the template]
Right now \texttt{my\_team.py} contains the same agents as \texttt{baseline\_team}, so games usually end in a \textbf{tie}.
That is expected. Results change once you write your own agents in \texttt{my\_team.py}.
\end{careful}

\begin{step}{Step 5 --- Fast games without graphics}
\begin{lstlisting}
python capture.py -r baseline_team -b ../../../my_team.py -q --delay-step 0
\end{lstlisting}
\begin{expect}
Within a second: \texttt{Red team starts} / \texttt{Time is up.} / the result line.
\end{expect}
\textbf{Why \texttt{--delay-step 0}?} By default the engine sleeps 0.03\,s after every move: 36\,s of waiting per game,
which looks like a hang. The tournament uses \texttt{--delay 0.0}, which does the same thing.
\end{step}

%=====================================================================
\section{All useful options}
%=====================================================================

Add any of these to the \texttt{python capture.py \ldots} command (still in \p{pacman-contest/src/contest}).

\begin{center}
\begin{tabularx}{\linewidth}{@{}lX@{}}
\toprule
\textbf{Option} & \textbf{What it does} \\
\midrule
\p{-l officeCapture} & Use a standard map. Choices: \texttt{defaultCapture} (default), \texttt{alleyCapture}, \texttt{bloxCapture},
  \texttt{crowdedCapture}, \texttt{distantCapture}, \texttt{fastCapture}, \texttt{jumboCapture}, \texttt{mediumCapture},
  \texttt{officeCapture}, \texttt{strategicCapture}, \texttt{testCapture}, \texttt{tinyCapture} \\
\p{-l RANDOM13} & Random map from seed 13 (same seed, same map). Any number works. \\
\p{-l RANDOM} & A new random map each game. \\
\p{-n 10} & Play 10 games in a row, then print \texttt{Scores}, win rates and \texttt{Record}. \\
\p{--delay-step 0.1} & Slow motion (seconds per move). \texttt{0} = full speed. \\
\p{-z 0.6} & Window size (smaller than 1 shrinks it). \\
\p{-q} & No graphics, minimal output. \\
\p{-Q} & No graphics, and hide your agents' own prints. \\
\p{-t} & Draw the game as text in Terminal. \\
\p{-i 400} & Shorter game: 400 total moves instead of 1200. \\
\p{-f} & Fixed random seed, so every run is identical (for debugging). \\
\p{--red-name A --blue-name B} & Team names shown in the window. \\
\p{--help} & List every option. \\
\bottomrule
\end{tabularx}
\end{center}

\subsection{Exact examples}
\begin{lstlisting}
# your team as RED instead of blue (the tournament plays both sides)
python capture.py -r ../../../my_team.py -b baseline_team

# office map, slow motion
python capture.py -r baseline_team -b ../../../my_team.py -l officeCapture --delay-step 0.1

# 10 fast games on random maps, with a summary at the end
python capture.py -r baseline_team -b ../../../my_team.py -l RANDOM -n 10 -q --delay-step 0

# your team against itself
python capture.py -r ../../../my_team.py -b ../../../my_team.py
\end{lstlisting}

\begin{careful}[Test both sides]
In the tournament your team plays as red \emph{and} as blue on random maps. Always test with \texttt{-r} and \texttt{-b} swapped, and with \texttt{-l RANDOM}.
\end{careful}

%=====================================================================
\section{Recording and replaying games}
%=====================================================================

\begin{step}{Record}
\begin{lstlisting}
python capture.py -r baseline_team -b ../../../my_team.py -q --delay-step 0 --record
\end{lstlisting}
Saved as \p{www/contest_default/replays/match_0.replay} (inside \p{pacman-contest/src/contest}).
Add \texttt{-m 7} to save it as \texttt{match\_7} instead, so the next recording doesn't overwrite it.
\end{step}

\begin{step}{Replay with graphics}
\begin{lstlisting}
python capture.py --replay www/contest_default/replays/match_0.replay --delay-step 0.05
\end{lstlisting}
For tournament replays downloaded from the results page:
\p{python capture.py --replay ~/Downloads/NAME.replay}
\end{step}

The \texttt{www/} folder only holds results and recordings. Delete it whenever you like: \texttt{rm -rf www}.

%=====================================================================
\section{Playing against your agent yourself}
%=====================================================================
\begin{lstlisting}
python capture.py -r baseline_team -b ../../../my_team.py --keys0
\end{lstlisting}
You control red agent 0. Click the game window first, then use the \textbf{arrow keys} or \textbf{W A S D} (\textbf{Q} = stop).

%=====================================================================
\section{Checking that it will work in the competition}
%=====================================================================

\texttt{CONTEST.md} says the organisers run every match like this:
\begin{lstlisting}
python3 capture.py --delay 0.0 -r teams/<team1>/my_team.py -b teams/<team2>/my_team.py
\end{lstlisting}
You can reproduce that exactly. From \p{pacman-contest/src/contest}:
\begin{lstlisting}
mkdir -p /tmp/teams/mine /tmp/teams/other
cp ../../../my_team.py /tmp/teams/mine/my_team.py
cp baseline_team.py   /tmp/teams/other/my_team.py
python3 capture.py --delay 0.0 -r /tmp/teams/mine/my_team.py -b /tmp/teams/other/my_team.py
\end{lstlisting}
If this game runs to the end without an error, your file will load in the tournament.

\textbf{Checklist before submitting}
\begin{itemize}[label=$\square$]
  \item The competition-style test above runs without errors, with your team as red and as blue.
  \item No \texttt{took too long to make a move} warnings (1\,s per move; the third warning loses the game).
  \item \texttt{create\_team()} works with its default arguments. The tournament passes no extra options.
  \item No \texttt{print(...)} left in \texttt{my\_team.py}. Everything printed is published with the results.
  \item Only Python's standard library plus \texttt{numpy}/\texttt{scipy}/\texttt{sklearn}/\texttt{tensorflow}/\texttt{keras}, which the organisers guarantee.
  \item \texttt{TEAM.md} filled in.
\end{itemize}

\begin{careful}[Where to submit]
The provided files do not say how this contest collects submissions. Follow the organisers' instructions for that step.
This folder's git \texttt{origin} is the organisers' template (\p{github.com/aig-upf/pacman-agent}), so you cannot push to it.
If they ask for a repository, create your own \textbf{private} GitHub repository and point the folder at it:
\begin{lstlisting}
git remote set-url origin git@github.com:<your-user>/pacman-agent.git
git add my_team.py TEAM.md && git commit -m "Submission"
git push -u origin main
\end{lstlisting}
\end{careful}

%=====================================================================
\section{Troubleshooting}\label{sec:trouble}
%=====================================================================

\begin{center}
\begin{tabularx}{\linewidth}{@{}>{\raggedright\arraybackslash}p{5.4cm}X@{}}
\toprule
\textbf{Message / symptom} & \textbf{Cause and exact fix} \\
\midrule
Window never appears, Python at 100\% CPU & Apple's old Python (Tk 8.5) is being used. Run \p{source venv/bin/activate}. If \texttt{TkVersion} still isn't 9.1, rebuild the venv (Section~\ref{sec:rebuild}). \\
\texttt{No module named 'contest'} or \texttt{'frozendict'} & venv not active, or packages missing. Activate it, then from \texttt{pacman-agent}: \p{pip install -e pacman-contest frozendict numpy}. \\
\texttt{No teams found!} / \texttt{could not be loaded} & You are not in \p{pacman-contest/src/contest} (Step~3), or the path to \texttt{my\_team.py} is wrong. \\
\texttt{The layout \ldots\ cannot be found} & Same cause: run from \p{pacman-contest/src/contest}. \\
Each game takes about 40 seconds & Add \p{--delay-step 0}. \\
\texttt{KeyError: 'game\_id'} & You ran \texttt{runner.py}. It is broken in the provided files. Use \texttt{python capture.py} as above. \\
\texttt{scripts/play.sh} gives \texttt{No teams found!} & It points to \texttt{myTeam.py}, which doesn't exist. Don't use it; use the commands above. \\
\texttt{replay.py}: \texttt{python3.8} not found & The provided \texttt{replay.py} is hard-wired to Python 3.8. Use \p{python capture.py --replay FILE}. \\
\texttt{Agent 1 took too long to make a move!} & Your agent needed more than 1\,s. Make it faster. \\
\texttt{zsh: command not found: python} & venv not active (Step~2). \\
\bottomrule
\end{tabularx}
\end{center}

\subsection{Rebuilding the environment from scratch}\label{sec:rebuild}
Only needed if the venv is broken or deleted:
\begin{lstlisting}
cd /Users/bortekstor/Desktop/Eutopia_Pacman_Contest/Pack_It_Up/pacman-agent
deactivate 2>/dev/null; rm -rf venv
/opt/homebrew/bin/python3.14 -m venv venv
source venv/bin/activate
pip install --upgrade pip
pip install -e pacman-contest
pip install frozendict numpy
python -c "import tkinter, contest; print(tkinter.TkVersion)"   # must print 9.1
\end{lstlisting}
If \texttt{/opt/homebrew/bin/python3.14} is missing: \p{brew install python@3.14 python-tk@3.14}.
Do \textbf{not} build the venv with plain \texttt{python3} (Apple's 3.9 with the broken Tk), and don't run \texttt{scripts/config.sh}
(it adds the engine from a different repository).

\subsection{Updating the game engine}
If the organisers announce an engine update, from \texttt{pacman-agent}:
\begin{lstlisting}
git submodule update --init --remote
source venv/bin/activate && pip install -e pacman-contest
\end{lstlisting}

\end{document}
